% !Mode:: "TeX:UTF-8"
% 北京工业大学数据挖掘课程报告格式定义

\makeatletter

% ============================================================
% 0. 兼容性定义
% ============================================================
\newcommand{\erhao}{\zihao{2}}
\newcommand{\xiaoer}{\zihao{-2}}
\newcommand{\sanhao}{\zihao{3}}
\newcommand{\xiaosan}{\zihao{-3}}
\newcommand{\sihao}{\zihao{4}}
\newcommand{\xiaosi}{\zihao{-4}}
\newcommand{\wuhao}{\zihao{5}}
\newcommand{\xiaowu}{\zihao{-5}}
\newcommand{\song}{\songti}
\newcommand{\hei}{\heiti}
\newcommand{\kai}{\kaishu}
\newcommand{\fs}{\fangsong}

% ============================================================
% 1. 基础行网格 (正文 36 行)
% ============================================================
\renewcommand\normalsize{
  % 重置行距倍数为 1.0，防止倍数叠加
  \renewcommand{\baselinestretch}{1.0}

  % 计算逻辑：版心 247mm / 36行 ≈ 19.45pt (微调值，防止溢出)
  \@setfontsize\normalsize{12pt}{19.45pt}

  % 强制取消段落间距
  \setlength{\parskip}{0pt}

  % 调整公式间距
  \setlength\abovedisplayskip{6pt plus 2pt minus 2pt}
  \setlength\belowdisplayskip{6pt plus 2pt minus 2pt}
  \setlength\abovedisplayshortskip{0pt plus 2pt}
  \setlength\belowdisplayshortskip{4pt plus 2pt minus 2pt}
  \let\@listi\@listI
}
\def\defaultfont{\normalsize\selectfont}
\setlength{\parindent}{2em}
\punctstyle{hangmobanjiao}
\setlist{nosep}

% ============================================================
% 2. 章节标题格式
% ============================================================
% 设置编号深度为 3 (显示 1.1.1.1)
\setcounter{secnumdepth}{3}

% 一级标题 (Chapter)
% 保持原样，或者段后也可以稍微改小
\titleformat{\chapter}{\centering\zihao{-3}\heiti}{\thechapter}{1em}{}
\titlespacing{\chapter}{0pt}{6pt}{6pt}

% 二级标题 (Section) - 1.1
% 将段后间距设为 0pt，防止标题高度超标导致跳行
\titleformat{\section}{\zihao{4}\heiti}{\thesection}{1em}{}
\titlespacing{\section}{0pt}{6pt}{0pt}

% 三级标题 (Subsection) - 1.1.1
% 为保持视觉一致，段后间距设为 0pt，段前间距设为 6pt
\titleformat{\subsection}{\zihao{4}\heiti}{\thesubsection}{1em}{}
\titlespacing{\subsection}{0pt}{6pt}{0pt}

% 四级标题 (Subsubsection) - 1.1.1.1
\titleformat{\subsubsection}{\zihao{-4}\heiti}{\thesubsubsection}{1em}{}
%\titlespacing{\subsubsection}{0pt}{0pt}{0pt}
\titleformat{\subsubsection}{\zihao{-4}\heiti}{\thesubsubsection}{1em}{}
\titlespacing{\subsubsection}{0pt}{6pt}{0pt}

% ============================================================
% 3. 页眉页脚 (奇偶页不同)
% ============================================================
\usepackage{fancyhdr}

% --- 1. 定义一个通用的页眉页脚配置命令 (供后面调用) ---
\newcommand{\bjutHeaderConfig}{
    \fancyhf{} % 清空
    \fancyhead[CO]{\songti\zihao{-5} 北京工业大学数据挖掘课程报告}
    \fancyhead[CE]{\songti\zihao{-5} \leftmark}
    \fancyfoot[C]{\rmfamily\zihao{-5} \thepage}
    \renewcommand{\headrulewidth}{0.5pt}
}

% --- 2. 设置默认的 fancy 样式 (用于正文非首页) ---
\pagestyle{fancy}
\bjutHeaderConfig  % <--- 关键：让普通页加载配置！

% --- 3. 设置默认的 plain 样式 (用于章首页) ---
% 默认情况下，我们在 format 里让 plain 和 fancy 长得一样
% 这样正文里的章首页就会有页眉
\fancypagestyle{plain}{
    \bjutHeaderConfig
}

% 章节标记格式
\renewcommand{\chaptermark}[1]{\markboth{\thechapter\quad #1}{}}
% ============================================================
% 4. 目录样式
% ============================================================
\renewcommand{\contentsname}{\zihao{3}\heiti 目\quad 录}
\setcounter{tocdepth}{2}

\newcommand{\worddots}{\titlerule*[3pt]{.}}

% --- 章 (一级标题) ---
% 格式：顶格 (无缩进要求)
\titlecontents{chapter}[0pt]
              {\vspace{0pt}\normalsize\songti\bfseries}
              {\rmfamily\thecontentslabel\hspace{0.5em}}
              {}
              {\hspace{0.5em}\worddots\contentspage}

% --- 节 (二级标题) ---
% 要求：之前有 1 个中文字符的空格 -> [1\ccwd]
\titlecontents{section}[1\ccwd]
              {\vspace{0pt}\normalsize\songti}
              {\rmfamily\thecontentslabel\hspace{0.5em}}
              {}
              {\hspace{0.5em}\worddots\contentspage}

% --- 小节 (三级标题) ---
% 要求：之前有 2 个中文字符的空格 -> [2\ccwd]
\titlecontents{subsection}[2\ccwd]
              {\vspace{0pt}\normalsize\songti}
              {\rmfamily\thecontentslabel\hspace{0.5em}}
              {}
              {\hspace{0.5em}\worddots\contentspage}
% ============================================================
% 5. 参考文献环境
% ============================================================
\renewenvironment{thebibliography}[1]{
    % 标题带编号
    \chapter{\bibname}

    % 直接使用正文默认设置，不强制单倍行距
    % 这样行距就会自动继承 \normalsize 的 19.52pt

    \list{\@biblabel{\@arabic\c@enumiv}}
         {\settowidth\labelwidth{\@biblabel{#1}}
          \leftmargin\labelwidth
          \advance\leftmargin\labelsep
          % 如果列表项之间太松散，可以保留下面这行；如果想完全和正文段落一致，就注释掉
          %\setlength{\itemsep}{0pt}
          \setlength{\itemsep}{0pt}   % 条目与条目之间的额外间距
          \setlength{\parsep}{0pt}    % 条目内段落之间的间距 (这个最关键！)
          \setlength{\parskip}{0pt}   % 段落前后的间距
          \setlength{\topsep}{0pt}    % 列表与上方正文的间距
          \setlength{\partopsep}{0pt} % 列表与上方页面的间距
          \usecounter{enumiv}
          \let\p@enumiv\@empty
          \renewcommand\theenumiv{\@arabic\c@enumiv}}
    \sloppy
    \clubpenalty4000
    \@clubpenalty \clubpenalty
    \widowpenalty4000
    \sfcode`\.\@m}
   {\def\@noitemerr
     {\@latex@warning{Empty `thebibliography' environment}}
    \endlist}
% ============================================================
% 6. 其他配置
% ============================================================
\renewcommand{\tablename}{表}
\renewcommand{\figurename}{图}
\renewcommand{\thefigure}{\thechapter-\arabic{figure}}
\renewcommand{\thetable}{\thechapter-\arabic{table}}
\renewcommand{\theequation}{\thechapter-\arabic{equation}}

\DeclareCaptionFont{captionwuhao}{\zihao{5}}
\DeclareCaptionFont{captionhei}{\heiti}
\captionsetup{font=captionwuhao, labelsep=quad}
\captionsetup[figure]{name=图, labelfont=captionhei, textfont=captionhei}
\captionsetup[table]{name=表, labelfont=captionhei, textfont=captionhei}

\def\temp{\relax}
\let\temp\addcontentsline
\gdef\addcontentsline{\phantomsection\temp}

% 项目成果列表
% \newlist{publist}{enumerate}{1}
% \setlist[publist]{label={[\arabic*]}, leftmargin=3em, itemsep=0pt, parsep=0pt, topsep=0pt}
\newlist{publist}{enumerate}{1}
\setlist[publist]{
    label={[\arabic*]},
    leftmargin=3em,
    itemsep=0pt,  % 条目间距
    parsep=0pt,   % 段落间距
    topsep=0pt,   % 与上方文本的间距
    partopsep=0pt,
    % 让列表内部行距跟随正文
    % 如果觉得太紧，可以试试 \setstretch{1.3} 之类的
    % 但既然要求“单倍行距”，这里默认跟随当前环境即可
}


% \def\cleardoublepage{\clearpage\if@twoside \ifodd\c@page\else
%   \hbox{}
%   \thispagestyle{empty}
%   \newpage
%   \if@twocolumn\hbox{}\newpage\fi\fi\fi}

% ============================================================
% 7. 覆盖 ctex 的默认编号设置
% ============================================================
\ctexset{
    chapter = {
        number = \arabic{chapter},
        name = {}, % 去掉“第”和“章”
    },
    section = {
        number = \thechapter.\arabic{section},
        name = {},
    },
    subsection = {
        % 关键修复：原本这里可能有 name={（,）}，我们要把它清空
        number = \thesection.\arabic{subsection},
        name = {},
        format = \zihao{4}\heiti, % 确保字体也是黑体
    },
    subsubsection = {
        number = \thesubsection.\arabic{subsubsection},
        name = {},
        format = \zihao{-4}\heiti,
    }
}
% ============================================================
% 8. 图表与公式格式 (修正双语图注前缀)
% ============================================================
\renewcommand{\tablename}{表}
\renewcommand{\figurename}{图}
\renewcommand{\thefigure}{\thechapter-\arabic{figure}}
\renewcommand{\thetable}{\thechapter-\arabic{table}}
\renewcommand{\theequation}{\thechapter-\arabic{equation}}

% 注册字体
\DeclareCaptionFont{captionwuhao}{\zihao{5}}
\DeclareCaptionFont{captionhei}{\heiti}
\DeclareCaptionFont{captioneng}{\rmfamily} % 英文用衬线体

% 基础设置
\captionsetup{
    font=captionwuhao,
    labelsep=quad, % 序号后空一格
    singlelinecheck=true % 单行居中
}

% --- 双语图注核心配置 ---
% lang=1 (中文) 配置
\captionsetup[figure][bi-first]{
    name=图,
    labelfont=captionhei,
    textfont=captionhei
}
% lang=2 (英文) 配置：关键是这里！
\captionsetup[figure][bi-second]{
    name=Fig.,           % 前缀改为 Fig.
    labelfont=captioneng,
    textfont=captioneng
}

% 表格同理
\captionsetup[table][bi-first]{
    name=表,
    labelfont=captionhei,
    textfont=captionhei
}
\captionsetup[table][bi-second]{
    name=Table,          % 前缀改为 Table
    labelfont=captioneng,
    textfont=captioneng
}

% ============================================================
% 9. 代码与算法的中文化重命名
% ============================================================
\renewcommand{\lstlistingname}{代码} % 将 Listing 改为 代码
\renewcommand{\algorithmcfname}{算法} % 将 Algorithm 改为 算法

% 修正 listings 的编号格式为 "章-序号" (如 代码 3-1)
\AtBeginDocument{
  \renewcommand{\thelstlisting}{\thechapter-\arabic{lstlisting}}
}

% ============================================================
% 10. 智能空白页控制 (由 metadata.tex 中的开关控制)
% ============================================================
\usepackage{ifthen} % 需要引入逻辑判断宏包

\makeatletter
\def\cleardoublepage{
  \clearpage
  % 如果是双面打印模式
  \if@twoside
    % 检查开关是否为 true
    \ifthenelse{\equal{\isOpenBlankPage}{true}}{
      % --- 模式 A：保留空白页 (标准做法) ---
      \ifodd\c@page
        % 当前是奇数页时，插入空白页，保证新章节从奇数页开始。
        \else
        % 当前是偶数页，不做操作，新章节直接从奇数页开始。
      \fi
      % 实际上 LaTeX 默认的 \cleardoublepage 就够了，
      % 但为了处理页眉页脚，我们通常会重写它：
      \ifodd\c@page\else
        \hbox{}
        \thispagestyle{empty} % 空白页不显页码
        \newpage
        \if@twocolumn\hbox{}\newpage\fi
      \fi
    }{
      % --- 模式 B：强制删除空白页 (紧凑做法) ---
      % 直接什么都不做，等于 \clearpage
      % 这样新章节就会紧接着上一章结束的页面开始
    }
  \fi
}

\bibpunct{[}{]}{,}{s}{}{;}

\makeatother
